




% !TeX program = pdflatex
% !BIB program = biber
\documentclass[11pt]{article}

\usepackage[a4paper,margin=1in]{geometry}
\usepackage[T1]{fontenc}
\usepackage[utf8]{inputenc}
\usepackage{lmodern}
\usepackage{amsmath,amssymb,amsthm,bm,mathtools}
\usepackage{booktabs,array,tabularx,longtable}
\usepackage{ragged2e}
\usepackage{enumitem}
\usepackage{xcolor}
\usepackage{xurl}
\usepackage[expansion=false]{microtype}
\setlength{\emergencystretch}{3em}
\usepackage{hyperref}
\usepackage[nameinlink,noabbrev]{cleveref}
\usepackage[backend=biber,style=numeric,sorting=none,maxnames=6,
            uniquename=false,url=false]{biblatex}
\addbibresource{cmotion_foundational_general.bib}
\addbibresource{cmotion_primary_research.bib}
\addbibresource{cmotion_internal_manuscripts.bib}


\AtEveryBibitem{\clearfield{note}\clearfield{annotation}\clearfield{keywords}}


\newtheorem{theorem}{Theorem}[section]
\newtheorem{proposition}[theorem]{Proposition}
\newtheorem{corollary}[theorem]{Corollary}
\theoremstyle{definition}
\newtheorem{definition}[theorem]{Definition}
\newtheorem{remark}[theorem]{Remark}

\newcommand{\cmotion}{\mbox{C-Motion}}
\newcommand{\ElEmEr}{\mbox{ElEmEr}}
\newcommand{\dd}{\mathrm{d}}
\newcommand{\kb}{k_{\mathrm B}}
\newcommand{\GammaFine}{\Gamma_{\mathrm{fine}}}
\newcommand{\GammaRed}{\Gamma_{\mathrm{red}}}
\newcommand{\GammaM}{\Gamma_{M}}
\newcommand{\LambdaC}{\Lambda_{C}}


\newcommand{\statusStructural}{\textcolor{blue!55!black}{\textbf{Structural}}}
\newcommand{\statusConditional}{\textcolor{yellow!60!black}{\textbf{Conditional}}}
\newcommand{\statusConjectural}{\textcolor{magenta!65!black}{\textbf{Conjectural}}}
\newcommand{\statusOpen}{\textcolor{red!75!black}{\textbf{Open}}}
\newcommand{\statusDeferred}{\textcolor{gray!70!black}{\textbf{Deferred}}}
\newcommand{\statusNegative}{\textcolor{red!65!black}{\textbf{Negative}}}


\newcommand{\natNOne}{\textcolor{blue!70!black}{\textbf{N1}}}
\newcommand{\natNTwo}{\textcolor{teal!70!black}{\textbf{N2}}}
\newcommand{\natControl}{\textcolor{gray!75!black}{\textbf{C}}}
\newcommand{\natOpen}{\textcolor{red!70!black}{\textbf{O}}}
\newcommand{\statusSharpOpen}{\textcolor{red!35!black}{\textbf{Open (sharp)}}}

\newcolumntype{Y}{>{\RaggedRight\arraybackslash}X}
\newcolumntype{P}[1]{>{\RaggedRight\arraybackslash}p{#1}}

\hypersetup{
  colorlinks=true,
  linkcolor=blue!50!black,
  citecolor=blue!50!black,
  urlcolor=blue!50!black,
  pdftitle={Thermodynamics in C-Motion: From Field States to Entropy, Temperature, and Equilibrium},
  pdfauthor={Toma Kavonius}
}


\title{\textbf{Thermodynamics in \cmotion{}:}\\
From Field States to Entropy, Temperature, and Equilibrium}
\author{Toma Kavonius}
\date{24 November 2026}

\begin{document}
\maketitle

\begin{center}
\small\itshape
Scope and limitations.  This manuscript develops a bounded thermodynamic interpretation
within C-Motion.  The native state-space measure, transition dynamics, equilibrium
selection, transport, quantum occupation, and microscopic arrow remain explicitly
unresolved where stated.
\end{center}
\vspace{0.6em}

\begin{abstract}
C-Motion takes resolved primitive propagation of its proposed Elemental
Motion Resonance (\ElEmEr{}) field to have local speed $c$.  Heating therefore
cannot mean that the primitive field moves faster.  We instead describe thermodynamic
variation through the organization of complete field states: recurrent closures,
open propagation, collective motion, internal excitation, spectral populations,
and correlations.  A macrostate discards selected fine-state distinctions;
entropy quantifies weighted multiplicity only after a physical reduction and
measure have been specified.

We separate the measure problem into physical state-space support, kinematic
reference weighting, dynamical invariance, and equilibrium weighting.  Neither
an invariant measure nor stationary coarse dynamics alone selects thermodynamic
equilibrium.  Temperature is the energy--entropy conjugate of the reduced
description, not a microscopic speed.  We distinguish dephasing, scattering,
energy exchange, occupation change, and full equilibration; repeated
closure-mediated interaction provides a conditional mixing route, not an
established microscopic thermalizer.

Heat and work are process-dependent accounts of transfer on a common,
conditionally conserved energy ledger; energy, action, and entropy remain
distinct.  Standard gas and radiation laws serve as downstream consistency
controls, with their kinetic, polarization, and occupation assumptions explicit.
Successor-directed causal propagation does not by itself establish a
thermodynamic arrow.  Native state-space measure selection, transition laws and
rates, closure spectra, ultraviolet occupation physics, transport, and robust
irreversibility remain open.
\end{abstract}


\section{Introduction: the fixed-\texorpdfstring{$c$}{c} thermodynamic puzzle}
\label{sec:introduction}
The primitive \ElEmEr{} propagation speed in \cmotion{} is fixed locally at
$c$ \cite{CmotionMainFramework2026}.  Thermodynamic variation must therefore
occur in the organization and statistical populations of complete field states,
not in a spread of primitive speeds.  Open modes, persistent closures, their
centre-of-momentum motion, internal excitations and correlations provide possible
degrees of freedom.  This does not by itself establish which of those degrees of
freedom are thermally active.

The field-first ontology also changes the microscopic counting problem.
The complete local field state is the interaction object, not a list of
independently identifiable primitive packets \cite{CmotionMainFramework2026}.
A physically reduced state space, coarse observables, and a justified measure
must precede statistical entropy.  The subsequent question---how the dynamics
redistributes energy and selects equilibrium---is distinct from the definition
of entropy or temperature.  The paper organizes those requirements as
\begin{equation}
\boxed{
\begin{aligned}
\text{complete field/closure state}
&\longrightarrow \text{coarse macrostate}
\longrightarrow \text{entropy} \\
&\longrightarrow \text{temperature}
\longrightarrow \text{mixing dynamics}
\longrightarrow \text{equilibrium}.
\end{aligned}}
\label{eq:reader-chain}
\end{equation}
The Main Framework supplies the local-$c$, complete-state and energy/action
premises.  Quantum Foundations owns the status of $J_0$, $J_0=\hbar$ and
$E=\hbar\omega$ used in later controls \cite{CmotionQuantumFoundations2026}.
No native microscopic thermalizer or equilibrium measure is derived here.

\section{From the physical fine state to a thermodynamic macrostate}
\label{sec:fine-state}
\subsection{What the fine physical state contains}

A \emph{fine physical state} is the most detailed physical state being admitted before
thermodynamic information is discarded.  Denote its state space by $\GammaFine$ and one
fine state by $\gamma_{\mathrm f}\in\GammaFine$.  In the present framework this is not a
list of independently persistent primitive particles.  It is the complete field/closure
configuration at the chosen stage of physical succession.  Here a ``complete''
configuration is complete relative to a specified admissible state representation:
its field and closure data, constraints, and correlations must suffice for the
proposed evolution in the domain being modelled.  Writing such a state does not
itself supply a preferred physical global simultaneity slice, a universally
aligned succession label $s$, or a native rule for comparing distant stages.
The geometric and dynamical compatibility of local state descriptions remains
part of the open microscopic evolution problem.

Depending on the physical regime, that configuration can include open propagating
modes; recurrent, localized organizations called \emph{closures}; the centre-of-momentum
motion of closures; internal closure organization and internal modes; correlations
between different regions or organizations; and extended interacting configurations.
An \emph{internal closure mode} simply means a recurrent degree of freedom within a
closure that can in principle carry or exchange energy without being identical to the
motion of the closure as a whole.

A simple reconstruction thought experiment shows why a scalar density is not enough to
specify this fine state.  Reproduce the same local energy density everywhere but change the
propagation directions, phases, or correlations.  The reconstructed state will in general
evolve differently.  Likewise, two configurations can have the same total momentum---even
$\bm P_{\rm tot}=0$---while one has strongly opposed directional structure and another is
nearly isotropic.  Equal values of a few coarse observables do not imply the same physical
microstate.

Phase, polarization, direction, spectrum, internal organization, and correlations are
therefore examples of fine-state distinctions.  They are not automatically separate entropy
``channels.''  Whether a distinction contributes to thermodynamic multiplicity depends on
what the macro-description resolves and on the physical weighting of the reduced state space.

This distinction matters because two states can have the same total energy and still be
physically organized very differently.  Schematically one may keep track of
\begin{equation}
  \mathcal E_{\rm open},\qquad
  \mathcal E_{\rm closure\,COM},\qquad
  \mathcal E_{\rm internal},\qquad
  \mathcal E_{\rm corr},
  \label{eq:organizational-ledger}
\end{equation}
without treating these labels as separately conserved forms of energy.  They are ways of
asking where the energy and state information reside within one physical configuration.

\subsection{When a closure coordinate is thermally active}
\label{sec:thermal-activity-gate}

An internal coordinate is not automatically a thermodynamic degree of freedom.
For a proposed closure mode, the following tests supply a
\emph{structural diagnostic}: (i) it must distinguish physical states
rather than coordinate or gauge descriptions; (ii) redundancy must be removed;
(iii) its variations must belong to the sector being counted rather than changing
an exactly fixed charge, topology or representation label; (iv) accessible
finite-energy excitations must exist without changing the closure species;
(v) the coordinate needs genuine dynamical/conjugate response; and
(vi) allowed interactions must connect the states whose thermal weights are
being assigned. These tests do not independently derive a mode or an
occupation probability.

Thus centre-of-momentum motion is an external contribution, not an internal
mode. Exact topological and action-like charges held fixed by a macrostate
condition the sector rather than supplying freely mixed degeneracy. Relative
phase, deformation/shape and branch imbalance are candidate internal modes
only if their excitation and response are physical. Common phase, spatial
orientation and director/frame texture require careful symmetry and
redundancy tests; open-mode dressing and correlations may contribute to the
full interacting state without being intrinsic isolated-closure levels.
Consequently,
\begin{equation}
  \boxed{\text{distinguishable closure coordinate}
  \;\not\Rightarrow\;\text{thermally accessible mode}.}
  \label{eq:closure-thermal-activity-guard}
\end{equation}

Even the frequencies of a physical small-deformation sector are not fixed
by its stiffness alone. As a \emph{standard quadratic control}, suppose
physical perturbation coordinates $q^a$ admit
\begin{equation}
  L^{(2)}=\frac12 G_{ab}\,q^{\prime a}q^{\prime b}
           -\frac12 H_{ab}\,q^a q^b,
  \qquad
  \det(H-\Omega^2G)=0.
  \label{eq:closure-spectrum-control}
\end{equation}
Here prime labels differentiation in the stipulated evolution description;
$H$ is the stiffness Hessian and $G$ the supplied nondegenerate kinetic
metric. This is not a claim that native \ElEmEr{} dynamics must be a
second-order quadratic system: a first-order/symplectic theory would require
its own spectral problem. A positive stiffness Hessian cannot by itself
supply the physical frequencies, their degeneracies or their thermal
occupation. Native kinetic/action structure, mode selection and accessible
excitation channels remain Open.

An intermediate-temperature regime illustrates the distinction. In a supplied
equilibrium model with internal gap $\Delta_{\rm int}$, a dilute gas can have
$T_{\rm deg}\ll T\ll\Delta_{\rm int}/\kb$: internal excitations are suppressed
while translational and configurational entropy can remain appreciable.
Here $T_{\rm deg}$ denotes the scale below which quantum degeneracy invalidates
the classical gas description; such a window need not exist for every model.
This is not a $T\to0$ claim. For the standard equilibrium quantum ideal-gas
control without extensive ground-state degeneracy, entropy tends to zero
as $T\to0$. No native C-Motion third law follows.
Slow or absent transitions imply kinetic arrest, not a third-law theorem;
exact disconnected sectors must not be pooled into one accessible ground-state
multiplicity without a justified physical mixing rule.

\subsection{A macrostate is a deliberate loss of detail}

Thermodynamics rarely tracks the fine state.  A \emph{macrostate} is a coarse description
that declares many fine states equivalent because they agree on a chosen set of
macroscopic observables.  We write this schematically as
\begin{equation}
  \GammaFine \xrightarrow{\ \mathcal C\ } \mathcal M,
  \qquad
  M=\mathcal C(\gamma_{\mathrm f}),
  \label{eq:macro-map-simple}
\end{equation}
where $\mathcal C$ is the coarse description (or macro-map) and $M$ is the resulting
macrostate.

Coarse-graining discards distinctions not resolved by specified observables,
controls and subsystem boundaries; it does not declare those distinctions
physically nonexistent.  States with the same pressure or energy, for example,
may still differ in phase, internal correlations or propagation structure.

For \cmotion{}, the difficult physical question lies upstream of the familiar entropy
formulas: which distinctions among open propagation, closures, internal modes,
correlations, and extended configurations should survive in the physical thermodynamic
state space?  The present framework has not yet derived a unique answer.

\subsection{Why primitive ray, crest, or packet counts are not entropy}

The field-first ontology rules out a tempting shortcut.  Historical or computational
labels used to describe a wave---a ray, crest, preparation packet, or propagation
segment---do not automatically remain independent physical identities after overlap,
interaction, or reorganization.  If the complete local field state is the physical
object, ancestry labels can cease to correspond to separately countable constituents.

Hence a schematic expression such as
\begin{equation}
  S\propto \kb\ln N_{\rm primitive}
  \label{eq:primitive-count-rejected}
\end{equation}
cannot be promoted to a physical entropy law merely because $N_{\rm primitive}$ can be
introduced as descriptive bookkeeping.

\begin{proposition}[No automatic primitive-count entropy]
\label{prop:no-count-entropy}
The field-first \cmotion{} ontology does not by itself supply a physically invariant,
independently countable primitive ray/crest/packet number from which thermodynamic
entropy can be defined.
\end{proposition}

This is a restriction, not a replacement entropy formula.  Counting can become
meaningful once the theory identifies physical reduced states or excitations that really
are distinct under the chosen thermodynamic description and supplies a rule for how
those states are weighted.

\section{Entropy begins with a physical reduction}
\label{sec:entropy-gate}
Entropy measures multiplicity or uncertainty only after the underlying alternatives have
been specified.  In a field-first theory the construction therefore has an explicit
physical order:
\begin{equation}
  \boxed{
  \GammaFine
  \longrightarrow
  \GammaRed
  \xrightarrow{\ \mathcal C\ }
  \mathcal M
  \longrightarrow
  \{\mu\ \hbox{or weights}\}
  \longrightarrow
  S.}
  \label{eq:entropy-pipeline}
\end{equation}
Here $\GammaRed$ is the physically reduced state space after genuine constraints and
redundancies have been handled; $\mathcal C$ maps a reduced state to a macrostate; and
$\mu$ is a physical measure or weighting rule.  The set of reduced states corresponding
to a macrostate $M$ is
\begin{equation}
  \Gamma_M=\{\gamma\in\GammaRed:\mathcal C(\gamma)=M\}.
  \label{eq:macroregion}
\end{equation}
A \emph{physical measure} here means the rule that says how much statistical weight is
assigned to different admissible reduced states.  It is not supplied merely by writing
down coordinates on state space.

Phase, polarization, internal organization, correlations, and environmental records may
all distinguish fine states.  Whether they remain thermodynamically distinct depends on
the physical reduction.  Conversely, two mathematical descriptions related only by a
true redundancy should not be counted twice.

\paragraph{Recurrence is not an entropy formula.}
A recurrent closure occupies a constrained class of configurations relative
to some specified ambient comparison space. Its organization, codimension,
or apparent geometric specialness does not by itself determine a lower
thermodynamic entropy. It can still have physical internal excitations,
orientations, sectors and correlations, with unknown statistical weights.
If an exactly recurrent subspace has zero measure under an ambient continuous
measure, substituting that zero into $\ln\mu(\Gamma_M)$ would not calculate
its sector entropy. One first needs a justified \emph{induced or conditional
measure on the physical constrained sector}, followed by the same fixed
reference-normalization convention used below. Neither recurrence nor
conditioning alone selects equilibrium weighting.

Under a coarse description, fine-state information can move into
correlations that the chosen macroscopic observables do not track, even when
the underlying evolution preserves the complete state:
\begin{equation}
  \boxed{
  \text{simple resolved organization}
  \longrightarrow
  \text{distributed fine correlations}.}
  \label{eq:correlation-redistribution}
\end{equation}
A coarse entropy can then change even if the complete state remains, in principle,
distinguishable.  That picture is an interpretation of coarse-graining, not a microscopic
second-law theorem; a monotonic entropy law still requires the appropriate dynamics and
boundary/statistical conditions.

\begin{proposition}[Entropy requires a physical state space]
\label{prop:reduction-gate}
A \cmotion{} entropy construction requires a physically admissible reduced state space,
a thermodynamic macro-map, and a physical measure or weighting rule.  These ingredients
have not yet been derived uniquely from the microscopic framework.
\end{proposition}

Standard Boltzmann and Gibbs entropies apply once the physically reduced
state space and weighting are supplied \cite{Boltzmann1877,Gibbs1902,Shannon1948Mathematical}.
For a continuous measure, fix a common, macrostate-independent reference unit
$\mu_*>0$ and write $\widehat\mu=\mu/\mu_*$ so the measure inside the
logarithm is dimensionless; for dimensionless counting measure set $\mu_*=1$.
For a macroregion of finite positive volume,
\begin{equation}
  S_{\rm B}(M)=\kb\ln\widehat\mu(\Gamma_M),
  \label{eq:boltzmann-control}
\end{equation}
and for a normalized probability measure $P$ absolutely continuous with respect
to $\widehat\mu$, let $p=\dd P/\dd\widehat\mu$ and $\int p\,\dd\widehat\mu=1$.
The Gibbs/Shannon form, when the integral exists, is
\begin{equation}
  S_{\rm G}[P;\widehat\mu]
  =-\kb\int_{\GammaRed}p(\gamma)\ln p(\gamma)\,
  \dd\widehat\mu(\gamma),
  \label{eq:gibbs-control}
\end{equation}
with the discrete sum used when appropriate
\cite{Boltzmann1877,Gibbs1902,Shannon1948Mathematical}.  A common fixed reference-unit change shifts these entropy expressions by
an additive constant; changes to relative state weights or energy-dependent
reference conventions can alter physical predictions.  Neither the native
weighting nor its cell/action normalization follows from the formulas.
The reference measure $\widehat\mu$ is distinct from the invariant
fine-state probability $\mu_Q$ used below.

\subsection{What can be fixed before the microscopic thermalizer is known}
\label{sec:prethermalizer-structure}

The statement ``the measure is unknown'' hides several distinct physical questions.  It is
useful to separate
\begin{equation}
  \boxed{
  \begin{array}{ll}
  \mathrm{M1}: & \text{physical state space/support: which states exist?}\\
  \mathrm{M2}: & \text{kinematic reference measure: how is that space geometrically weighted?}\\
  \mathrm{M3}: & \text{dynamically invariant measure: which weighting is preserved?}\\
  \mathrm{M4}: & \text{thermodynamic equilibrium weighting: which weighting is physically realized?}
  \end{array}}
  \label{eq:measure-layers}
\end{equation}
The complete answers to M1, M3, and M4 remain open in \cmotion{}.
Some aspects of M2 can be determined from the symmetries assumed in the
effective models.  However, a measure preserved by the dynamics is not
necessarily the distribution that a physical system reaches in thermal equilibrium.

\begin{proposition}[Invariance alone does not select thermodynamic weighting]
\label{prop:invariant-measure-no-go}
For a deterministic evolution map $F$ with an invariant probability
$F_*\mu=\mu$, the invariance equation and conserved-sector labels do not
\emph{in general} identify a unique physically selected thermodynamic
weighting. This does not assert that every evolution admits multiple invariant
probability measures.
\end{proposition}
Fixed points, periodic orbits and mixtures of invariant measures provide
conditional counterexamples to automatic selection.  Uniquely ergodic systems
also exist, but uniqueness of an invariant probability does not alone identify
it with physical thermal equilibrium.  Hence
\begin{equation}
  \boxed{
  \text{dynamical invariance + conservation}
  \not\Rightarrow
  \text{unique thermodynamic measure}.}
  \label{eq:invariant-measure-no-go}
\end{equation}
Even a future Hamiltonian realization with a preserved Liouville-type volume would therefore
leave a separate physical question: why that weighting, rather than another invariant
component or mixture, is the thermodynamic one.

There is nevertheless a useful conditional kinematic result for closure centre-of-momentum
motion.  In the already-Conditional effective massive-closure/Lorentz sector, write the
operational four-momentum as $p^\mu=(p^0,\bm P)$, with $p^0=E/c$
and metric signature $(+---)$ as in the Main Framework, so that
\begin{equation}
  p^\mu p_\mu=M^2c^2,\qquad p^0>0.
  \label{eq:mass-shell-reference}
\end{equation}
The positive mass shell carries, up to normalization, the invariant reference element
\begin{equation}
  \dd\nu_M\propto\frac{\dd^3P}{p^0},
  \qquad
  \dd\nu_{\rm dir}=\frac{\dd\Omega}{4\pi}
  \label{eq:com-reference-measures}
\end{equation}
for the momentum shell and isotropic direction respectively, while homogeneity supplies the
ordinary physical-three-space element $\dd^3x$.  The four-component momentum notation is an
effective Lorentz bookkeeping device; it does not change the framework's primitive spatial
ontology from physical three-space.  The invariant element $\dd\nu_M$ and the
rest-frame coordinate element $\dd^3P$ are not the same reference measure.  For any
probability distribution described by densities $\rho_\nu$ relative to $\dd\nu_M$
and $f_3$ relative to $\dd^3P$, their common probability element obeys, up to the
normalization convention for $\dd\nu_M$,
\begin{equation}
  \rho_\nu(\bm P)\,\dd\nu_M
  = f_3(\bm P)\,\dd^3P,
  \qquad \rho_\nu(\bm P)\propto p^0 f_3(\bm P).
  \label{eq:com-density-conversion}
\end{equation}
This is a change-of-reference-measure identity, not an equilibrium law.  In particular,
the bare density $e^{-\beta E}$ relative to $\dd\nu_M$ would carry an extra
$1/p^0$ factor when expressed relative to $\dd^3P$.  Which density is physically
selected belongs to the separate statistical/equilibrium problem (M4).
Most importantly,
\begin{equation}
  \boxed{
  \text{kinematic reference measure}
  \neq
  \text{thermodynamic occupation weighting}.}
  \label{eq:reference-not-thermal-weight}
\end{equation}
The Boltzmann factor or any other equilibrium weighting still requires the statistical and
physical equilibrium argument.  Nor is a factorization of centre-of-momentum and internal
structure assumed.  Schematically one may instead disintegrate a joint measure as
\begin{equation}
  \dd\mu(P,Y)=\dd\mu_{\rm COM}(P)\,\dd\mu(Y\mid P),
  \label{eq:measure-disintegration}
\end{equation}
where $Y$ collects internal closure structure, dressing, and correlations.  The COM kinematic
geometry is partly constrained in the conditional sector; the internal/correlation weighting
remains open.

A second structural result concerns occupation variables.  Let $N_i$ be coarse occupations
and let $\nu_i^{(r)}$ denote an allowed reaction direction.  Chemical equilibrium at fixed
exact conserved quantities requires
\begin{equation}
  \sum_i\nu_i^{(r)}\mu_i=0
  \qquad\text{for every allowed reaction }r.
  \label{eq:chemical-equilibrium-network}
\end{equation}
If the exact conserved charges are
\begin{equation}
  Q_\alpha=\sum_i q_{\alpha i}N_i
  \label{eq:conserved-occupation-charges}
\end{equation}
and the charge map is $q:\mathbb R^{N}\to\mathbb R^{m}$,
$(qv)_\alpha=\sum_iq_{\alpha i}v_i$, assume precisely
$\operatorname{span}\{\nu^{(r)}\}=\ker q$ in the supplied finite occupation
space. At an interior equilibrium where the stated reaction variations and
chemical potentials are defined, $\mu$ annihilates $\ker q$ and hence belongs
to $(\ker q)^\perp=\operatorname{im}q^{\mathsf T}$:
\begin{equation}
  \boxed{
  \mu_i=\sum_\alpha\lambda_\alpha q_{\alpha i}.}
  \label{eq:chemical-potential-charge-span}
\end{equation}
Thus equilibrium chemical potentials lie in the span of the exact conserved
charge vectors under this full-span hypothesis. For a resolved open-mode
occupation $N_{\rm o}$, it follows that $\mu_{\rm o}=0$ if
$q_{\alpha{\rm o}}=0$ for every $\alpha$. This is a sufficient implication,
not a necessary condition: charge multipliers may also cancel or vanish.
Occupation-changing processes alone do not establish the full-span hypothesis.
If an exact number or action-like charge protects that occupation, zero
chemical potential is not generally forced.  This is an equilibrium constraint on a
supplied reaction network, not a derivation of the native occupation law or of the required
conversion dynamics.

The same fine/coarse structure also sharpens what stochastic thermodynamics can mean in a
deterministic field theory.  This is a conditional dynamical-systems construction:
assume a compatible complete-state representation and a well-defined evolution
update within it.  The update label $\Delta$ denotes a specified finite step in
that representation; it is not thereby a universal $s$ increment, an observable
clock reading, or a fourth physical coordinate.  The construction does not
establish that a single global equal-$s$ state description exists in native
\cmotion{}.  Fix exact conserved quantities $Q$, let $\Gamma_Q$ be the
corresponding fine-state sector, let
\begin{equation}
  F_\Delta:\Gamma_Q\rightarrow\Gamma_Q
  \label{eq:fine-map}
\end{equation}
be deterministic complete-state evolution over the stipulated coarse step, and let
$C:\Gamma_Q\rightarrow\mathcal A$ define coarse cells $A_a=C^{-1}(a)$.  For this conditional construction, supply a normalized fine-state probability measure
$\mu_Q(\Gamma_Q)=1$.  Let the $A_a$ form a measurable partition (up to null sets), and
restrict the transition formula to cells with $\mu_Q(A_a)>0$.  The one-step conditional
coarse transition weights are then
\begin{equation}
  \boxed{
  W_{ab}
  =\frac{\mu_Q\!\left(A_a\cap F_\Delta^{-1}A_b\right)}{\mu_Q(A_a)}}
  \label{eq:coarse-one-step-kernel}
\end{equation}
with $W_{ab}\geq0$ and $\sum_bW_{ab}=1$.  The formula samples unresolved fine
states using $\mu_Q$ conditioned on the starting cell; it does not derive that measure
from native dynamics.  Such one-step stochasticity can arise even though the fine evolution
is deterministic.  The randomness reflects unresolved fine-state distinctions rather than
inserted microscopic noise.

One-step stochasticity does not, however, imply an autonomous Markov process.  If the fine
distribution at step $n$ is $\rho_n$, the exact projected transition can depend on the hidden
conditional distribution inside $A_a$,
\begin{equation}
  W_{ab}^{(n)}
  =\int_{A_a}\mathbf 1_{A_b}\!\left(F_\Delta(x)\right)\,\rho_n(\dd x\mid A_a).
  \label{eq:history-dependent-coarse-kernel}
\end{equation}
That conditional distribution can retain preparation history.  Hence
\begin{equation}
  \boxed{
  \text{coarse stochasticity}\neq\text{coarse Markovianity}.}
  \label{eq:stochastic-not-markov}
\end{equation}
An \emph{autonomous Markov description of successive coarse histories},
whose multi-step statistics are generated by one kernel, needs an additional
justification. Exact lumpability is one sufficient condition; for a supplied
finite fine-state Markov chain, the Kemeny--Snell criterion requires the same
total transition probability into each target cell from every fine state in
the same starting cell \cite{KemenySnell1960}. This finite-chain control is
not a proof of lumpability for the native continuum field evolution. Rapid
conditional mixing, timescale separation or bath resetting may support an
approximation only under further stated conditions. The fixed fine map,
partition and supplied fine-state probability $\mu_Q$ already specify the time-independent
\emph{one-step} matrix in \cref{eq:coarse-one-step-kernel}, even when the
actual multi-step histories retain memory. Whether physical closure
dynamics supports Markov closure remains Open.

If the normalized $\mu_Q$ is invariant under $F_\Delta$, then
\begin{equation}
  \pi_a=\mu_Q(A_a)
  \label{eq:stationary-pushforward}
\end{equation}
is a probability vector satisfying the one-step identity
\begin{equation}
  \boxed{\pi W=\pi,}
  \label{eq:global-stationarity}
\end{equation}
because $\sum_a\pi_aW_{ab}=\mu_Q(F_\Delta^{-1}A_b)=\mu_Q(A_b)$.
This identity requires no autonomous Markov closure: the invariant fine ensemble has
stationary coarse one-time marginals even when its multi-step coarse histories are
non-Markovian.  Only if an autonomous kernel is additionally justified may successive
coarse distributions generally be propagated as $P_{n+1}=P_nW$.  Pairwise detailed
balance is not needed for the one-step stationarity identity.  Defining
\begin{equation}
  J_{ab}=\pi_aW_{ab}-\pi_bW_{ba},
  \label{eq:coarse-probability-current}
\end{equation}
stationarity requires only zero net divergence, $\sum_bJ_{ab}=0$, not $J_{ab}=0$ on every
edge.  A stationary distribution may therefore carry circulating probability currents,
as in the cycle description of master-equation networks
\cite{Schnakenberg1976}; the displayed $J_{ab}$ is its discrete-step counterpart.
Pairwise detailed balance imposes the stronger condition
\begin{equation}
  \pi_aW_{ab}=\pi_bW_{ba}.
  \label{eq:detailed-balance-control}
\end{equation}
Consequently,
\begin{equation}
  \boxed{
  \text{preparation independence + stationarity}
  \not\Rightarrow
  \text{thermodynamic equilibrium}.}
  \label{eq:stationary-not-equilibrium}
\end{equation}
The positive physical criterion distinguishing equilibrium from a nonequilibrium stationary
state remains unresolved.

For an autonomous Markov kernel with stationary $\pi$, the standard
data-processing inequality \cite{KullbackLeibler1951,CoverThomas2006} gives
\begin{equation}
  D(PW\Vert\pi)\leq D(P\Vert\pi).
  \label{eq:relative-entropy-contraction}
\end{equation}
Thus relative entropy is nonincreasing for distributions with finite defined $D$;
it does not alone prove strict decrease or convergence $P_n\to\pi$.  For example,
a deterministic three-state cycle has uniform stationary distribution
$\pi=(1/3,1/3,1/3)$; starting at one state, the probability rotates forever while
$D(P_n\Vert\pi)=\ln 3$ for every $n$.  Finite-state irreducibility together with
aperiodicity is one sufficient additional condition for convergence, but is not obtained
from stationarity alone.  Nor does this inequality prove that $\pi$ is Gibbs, that
$\pi$ is thermodynamic equilibrium, that $D$ is the native thermodynamic entropy, or that
a second law has been derived.  Likewise, conservation, an allowed transition graph, and
connectivity do not select Gibbs weighting: even formal detailed-balance kernels can be built
around different positive stationary distributions.  The energy dependence of the physical
equilibrium weight must come from the physical measure/rate structure rather than from
connectivity or balance terminology alone.

These results narrow the microscopic problem without closing it.  The complete physical
state space, internal/correlation measure, selected invariant/equilibrium weighting, actual
transition law and rates, coarse memory-loss mechanism, physical connectivity, and positive
equilibrium-selection criterion remain to be supplied by the underlying dynamics.

\section{Temperature when primitive propagation is fixed}
\label{sec:temperature}
Return now to the motivating puzzle.  Fixing the characteristic propagation speed does not
fix the distribution of energy.  Even a purely open-mode control can have the same dispersion
speed $c$ with different spectral occupation functions; the two preparations then carry
different energies and can belong to different thermodynamic states.  With closures present,
the possibilities multiply: energy can reside in centre-of-momentum motion, internal
excitation, open modes, correlations, or different accessible organizations.  One system can
therefore be hotter than another without any primitive component ``speeding up.''

\begin{proposition}[Primitive local $c$ does not determine temperature]
\label{prop:c-not-temperature}
Fixing the local characteristic speed of primitive \ElEmEr{} propagation does not fix a
unique thermodynamic temperature.
\end{proposition}

The permitted temperature dependence must instead involve energy distribution,
mode occupation, composite motion or correlations.  It cannot generally be
identified with primitive speed, crest/ambient density or net momentum.

Once an entropy function and internal energy $U$ have been defined for the chosen
thermodynamic description, the usual conjugate definition applies:
\begin{equation}
  \boxed{
  \frac{1}{T}
  =\left(\frac{\partial S}{\partial U}\right)_{\{\mathcal X\}}.}
  \label{eq:temperature-control}
\end{equation}
The variables $\{\mathcal X\}$ denote whatever other macroscopic constraints are held
fixed.  The equation gives temperature its thermodynamic meaning: it describes how the
entropy of the accessible macrostate structure changes when the system's energy changes.
It is not a derivation of $T$ from the universal propagation speed.

Populating closure modes, changing collective motion or redistributing spectral
weight can change this derivative.  Which processes occur, and whether a thermal
contact regime exists, depend on the as-yet-unselected microscopic dynamics.

\paragraph{Operational thermometry.}
A derived $T$ should not remain just a formal derivative or fitted width.
In a specified equilibrium/contact regime, independently calibrated physical
probes weakly coupled to the same system should agree on one temperature
without a probe-dependent redefinition of the state measure. The native
interaction and equilibration law needed to establish that operational
agreement are still Open.

\subsection{Massive closure motion is compatible with fixed primitive speed}

A recurrent closure is an organized field pattern, not a single primitive ray.  Its
centre of momentum can move subluminally even though the primitive propagation that
sustains the pattern remains locally at $c$.  The closure may also possess internal modes
whose occupations change with energy.  This separates three notions that are easily
confused: primitive propagation, motion of the organized object as a whole, and internal
excitation of that object.

Consequently the ordinary kinetic intuition can reappear at the \emph{closure} level.
If a stable closure sector is supplied with the timelike dispersion
\begin{equation}
  E(\bm P)=\sqrt{M^2c^4+P^2c^2},
  \label{eq:timelike-control}
\end{equation}
then the conditional mass-shell geometry discussed in \cref{sec:prethermalizer-structure}
constrains a kinematic reference measure, but it does not supply the thermal occupation.
Once a standard \emph{rest-frame} equilibrium prescription is additionally supplied,
the familiar classical J\"uttner density \cite{Juttner1911} with respect to the coordinate element
$\dd^3P$ is
\begin{equation}
  f_3(\bm P)\propto
  \exp\!\left[-\frac{\sqrt{M^2c^4+P^2c^2}}{\kb T}\right].
  \label{eq:juttner-control}
\end{equation}
This is a density relative to $\dd^3P$, not relative to the invariant element
$\dd\nu_M\propto\dd^3P/p^0$.  Describing the same distribution relative to
$\dd\nu_M$ requires the $p^0$ conversion factor in
\cref{eq:com-density-conversion}; postulating a bare Boltzmann density relative to
$\dd\nu_M$ would define a different probability distribution.  The choice between
such prescriptions is not fixed by kinematic mass-shell invariance.
Cubero et al. compared the original and modified J\"uttner prescriptions
using relativistic simulations of a confined one-dimensional two-component
gas; their equal-laboratory-time velocity samples favored the original
J\"uttner distribution \cite{Cubero2007RelativisticThermometers}.
This is dynamical evidence in that specified model, not a kinematic selection
theorem or a universal C-Motion equilibrium measure. Their confined setup
also makes explicit that boundary and sampling conditions belong to the
thermometry problem; it does not supply native closure thermometry.

For an individual momentum $P\ll Mc$, the energy expands as
$E(\bm P)=Mc^2+P^2/(2M)+O(P^4/(M^3c^2))$.  For that approximation to describe
the bulk of a thermal ensemble in its rest frame, additionally assume a
nondegenerate classical regime and the nonrelativistic thermal condition
$\kb T\ll Mc^2$ (so typical populated momenta satisfy $P\ll Mc$).  After
absorbing the common rest-energy factor into normalization, the distribution
reduces to the Maxwell--Boltzmann momentum density
\begin{equation}
  f_{\rm MB}(\bm P)\propto
  \exp\!\left[-\frac{P^2}{2M\kb T}\right].
  \label{eq:maxwell-control}
\end{equation}
These distributions demonstrate compatibility, not a microscopic \cmotion{} derivation.  The
closure dispersion, internal spectrum, and equilibrium measure remain upstream inputs.  Nor
does every single isolated closure need to possess a sharp thermodynamic temperature.  That
concept is naturally an ensemble, bath, or sufficiently dense internal-state description;
a small closure with only a few resolved levels may instead be described by explicit state
populations.

\section{Thermalization is a process problem}
\label{sec:hierarchy}
Entropy and temperature definitions do not supply equilibration dynamics.
Dephasing can suppress coarse interference without moving energy; directional
scattering can produce isotropy without changing occupation; energy exchange
can leave protected particle or mode numbers fixed.  The physical thermalizer
must supply the relevant transitions and a justified equilibrium criterion.
Preparation loss alone is insufficient, as the stationary-state controls in \cref{sec:prethermalizer-structure} show.  The different dynamical jobs can be summarized by the following hierarchy.  The labels
are shorthand introduced only after the physical distinctions are clear:
\begin{center}
\renewcommand{\arraystretch}{1.24}
\begin{tabularx}{\textwidth}{@{}P{0.08\textwidth}P{0.29\textwidth}Y@{}}
\toprule
Level & Physical process & What it accomplishes \\
\midrule
T0 & Free propagation
& Carries the existing mode and correlation structure forward; no generic equilibration follows. \\
T1 & Dephasing / correlation redistribution
& Coarse coherences may relax while information survives in finer phase/correlation structure. \\
T2 & Angular / directional scattering
& Redistributes directions and can support isotropization. \\
T3 & Spectral or energy exchange
& Redistributes energy among frequencies, momenta, or internal levels. \\
T4 & Occupation-changing processes where required
& Changes the number or population of effective excitations when equilibrium requires such channels. \\
T5 & Repeated connected mixing
& Can erase preparation dependence only if the physical transition network is sufficiently connected; coarse closure, stationary weighting, and thermodynamic equilibrium identification are additional requirements. \\
\bottomrule
\end{tabularx}
\end{center}

Two negative statements follow immediately:
\begin{equation}
  \boxed{\text{dephasing}\neq\text{thermalization}},
  \qquad
  \boxed{\text{isotropy}\neq\text{full equilibrium}}.
  \label{eq:dephasing-negative}
\end{equation}
Free local-$c$ propagation can therefore contribute to T0 and, through phase mixing, to
some T1 behaviour without supplying the stronger redistribution required at T2--T5.  In
particular,
\begin{equation}
  \boxed{
  \text{loss of visible coherence}
  \not\Rightarrow
  \text{energy redistribution}
  \not\Rightarrow
  \text{full equilibrium}.}
  \label{eq:thermalization-ladder-negative}
\end{equation}
The hierarchy itself is generic.  The \cmotion{} question is which physically admissible
field and closure processes can actually perform each job.  The final step must now be read
with a sharper separation:
\begin{equation}
  \boxed{
  \text{physical connectivity}
  \neq
  \text{coarse Markov closure}
  \neq
  \text{stationarity}
  \neq
  \text{thermodynamic equilibrium}.}
  \label{eq:t5-sharpened}
\end{equation}
Connectivity specifies permitted transitions; coarse Markov closure requires
loss of relevant hidden memory; stationarity concerns a preserved distribution;
equilibrium needs an additional physically selected thermodynamic criterion.

\section{How closures could mediate equilibration}
\label{sec:matter-mediated}
One physically economical possibility is that freely propagating vacuum modes remain nearly
transparent to one another while organized matter supplies most of the redistribution.  An
incoming open mode can interact with a closure, transfer energy into collective or internal
organization, and later return energy to one or more outgoing modes.  A representative
sequence is
\begin{equation}
  \boxed{
  \begin{aligned}
  \text{open field mode}
  &\longrightarrow \text{closure interaction} \\
  &\longrightarrow \text{absorption / internal mixing / redirection} \\
  &\longrightarrow \text{re-release or further interaction}.
  \end{aligned}}
  \label{eq:closure-cycle}
\end{equation}

Repeated closure encounters could redistribute direction, phase, polarization,
frequency and, where allowed, mode occupations.  The proposed physical channels
must nevertheless be demonstrated, not inferred from the number of internal modes.

A useful schematic is
\begin{equation}
  \boxed{
  \text{open mode}+C_a
  \;\rightleftarrows\;
  C_b
  \;\rightleftarrows\;
  C_c+\text{open mode(s)},}
  \label{eq:closure-microcycle}
\end{equation}
where $C_a,C_b,C_c$ denote different closure microstates.  The diagram is only a process
picture: it assumes no derived transition kernel or rate law.

This closure-mediated mechanism remains conditional.  It identifies a plausible class of
physical processes capable in principle of supplying angular scattering, spectral exchange,
occupation change, and repeated mixing.  The actual transition map or generator, rates,
physical connectivity, intra-cell memory loss or other coarse-closure mechanism, and the
criterion selecting a thermodynamic equilibrium weighting are still unknown.  A large number
of internal degrees of freedom does not by itself turn a closure into a memoryless random
bath.

Fine information need not be destroyed in such encounters: correlations with
the closure and outgoing field can retain it even when the coarse subsystem
appears relaxed.  This is an interpretation consistent with the proposed
complete-field bookkeeping, not a derived thermalization law.

\subsection{How transparent free propagation could coexist with thermal contact}

The Main Framework allows a candidate picture in which ordinary free propagation can remain
largely unaffected until a local configuration departs from the free-propagation class and a
physical response channel becomes available \cite{CmotionMainFramework2026}.  In thermodynamic
language, mutually compatible open propagation could therefore remain comparatively
transparent, while encounters with recurrent or sufficiently reorganized field structure
could provide opportunities for redirection and energy exchange.

Such a picture would explain how an apparently transparent vacuum can coexist with efficient
thermalization inside matter.  It remains conjectural.  The framework does not show that
every non-null, non-collinear, or overlapping configuration interacts; it does not derive the
rates; and it does not establish that the same microscopic response forms closures, thermalizes
matter, or supplies a universal force.

\section{Energy transfer: heat, work, action, and entropy}
\label{sec:bookkeeping}
Thermodynamics does not introduce a second kind of physical energy.  The internal energy
$U$ of a thermodynamic subsystem is booked on the same physical energy ledger used by the
underlying framework.  The Main Framework separately distinguishes energy from action:
an independently conserved action charge requires the corresponding exact phase/action
symmetry \cite{CmotionMainFramework2026}.  The normalization slot $J_0$ belongs to that
action layer.  Its empirical identification with $\hbar$ is a Quantum-Foundations issue,
not a thermodynamic derivation \cite{CmotionQuantumFoundations2026}.

This common energy ledger is a \emph{conditional conservation and exchange premise},
not a microscopic energy-conservation theorem already derived from the native \ElEmEr{}
field law.  For a supplied compatible complete-state evolution $F_\Delta$ and an
isolated complete field-plus-environment state $\gamma$, assume there exists a
well-defined total-energy functional $\mathcal E$ such that
\begin{equation}
  \mathcal E(F_\Delta\gamma)=\mathcal E(\gamma)
  \qquad\text{(isolated-system energy control).}
  \label{eq:energy-conservation-control}
\end{equation}
The subsystem energy $U$ and its changes must then be defined by a specified boundary
and an accounting of the corresponding energy transfers.  Neither the full native
functional $\mathcal E$ nor its exact local transfer law has yet been supplied;
$F_\Delta$ remains a stipulated evolution in a compatible state representation, not a
globally established succession clock.  This premise is distinct from conservation
of the action charge associated with a separate symmetry.

Entropy belongs to a third layer.  It characterizes the multiplicity or uncertainty of a
reduced thermodynamic description.  It is therefore neither energy nor action, and an
action normalization cannot simply be reinterpreted as an entropy measure.

Under the supplied energy-conservation and boundary-accounting premise, with all
energy exchanges grouped into the selected heat and work channels, the first law
for the sign convention used here is
\begin{equation}
  \boxed{\dd U=\delta Q+\delta W.}
  \label{eq:first-law}
\end{equation}
If the chosen macroscopic description retains additional transfer channels separately,
the corresponding terms must be included.  The displayed two-term form is standard
bookkeeping, not a derivation of the native energy functional.  Its physical reading
is the important part.\ \emph{Work}
classifies energy transfer associated with controlled generalized changes of the chosen
system.\ \emph{Heat} classifies energy transfer through the remaining thermodynamic
interaction channels under that same coarse description.  The split depends on the
system boundary, the controls, and the macroscopic variables being tracked; heat and
work are not different substances stored inside the system.

Controlled boundary motion or coherent driving is work-like; unresolved
closure-mode redistribution is heat-like under a suitable subsystem reduction.
The classification depends on the controls and boundaries, not on an intrinsic
``heat'' label attached to an open field mode.

For the same reason,
\begin{equation}
  \boxed{\text{radiated energy}\not\equiv\text{heat}.}
  \label{eq:radiation-not-heat}
\end{equation}
A coherent outgoing field can remain work-capable, and energy released by an interaction
can remain correlated and nonthermal.  Radiation becomes heat-like only relative to a
specified thermodynamic process and reduction.

\section{Equilibrium and the downstream statistical description}
\label{sec:controls}
A stationary coarse distribution is not automatically thermal equilibrium;
it may retain currents or depend on unmodelled history.  A native equilibrium
claim must separately supply physical weighting, dynamically accessible
transitions and a criterion distinguishing thermal equilibrium from another
stationary state.

Once the physical state space, reference measure, conserved quantities, and
equilibrium antecedents have been specified, familiar statistical mechanics
provides the downstream language.  For example, maximizing the Gibbs entropy
of \cref{eq:gibbs-control} at fixed mean energy, with the same fixed
\emph{dimensionless} reference measure $\widehat\mu$, gives a density relative
to that measure:
\begin{equation}
  p_\beta(\gamma)
  =\frac{1}{Z(\beta)}\exp[-\beta H(\gamma)],
  \qquad
  Z(\beta)=\int\exp[-\beta H] \,\dd\widehat\mu,
  \qquad
  \beta=\frac{1}{\kb T},
  \label{eq:canonical-control}
\end{equation}
under the usual assumptions and for finite $Z$ \cite{Jaynes1957}.  The calculation does not determine the
microscopic \cmotion{} Hamiltonian, reduced state space, physical measure, or equilibrium
selection rule; those remain upstream.


Consider two weakly coupled systems with
\begin{equation}
  E_A+E_B=E_{\rm tot},
\end{equation}
and supplied physical multiplicities $\Omega_A(E_A)$ and $\Omega_B(E_B)$,
\begin{equation}
  P(E_A)\propto\Omega_A(E_A)\Omega_B(E_{\rm tot}-E_A).
  \label{eq:two-bath-weight}
\end{equation}
Writing $S=\kb\ln\Omega$, the most probable energy partition satisfies
\begin{equation}
  \frac{1}{T_A}=\frac{1}{T_B},
  \qquad\text{hence}\qquad
  \boxed{T_A=T_B}.
  \label{eq:equal-temperature-control}
\end{equation}
For positive heat capacities,
\begin{equation}
  \frac{\dd^2S_{\rm tot}}{\dd E_A^2}
  =-\frac{1}{T_A^2C_A}-\frac{1}{T_B^2C_B}<0,
  \label{eq:two-bath-stability}
\end{equation}
so the equal-temperature partition is a local entropy maximum.  For a small subsystem in a
large bath, the same multiplicity argument gives the standard canonical or grand-canonical
weighting
\begin{equation}
  P_a\propto
  \exp\!\left[-\beta\left(E_a-\sum_r\mu_rQ_{r,a}\right)\right].
  \label{eq:bath-canonical-control}
\end{equation}
These are standard equilibrium controls, conditional on a physical state measure;
they do not yield native transition rates or measure selection.

The stationarity, current and data-processing controls are established in
\cref{sec:prethermalizer-structure}; see
\cref{eq:global-stationarity,eq:coarse-probability-current,eq:detailed-balance-control,eq:relative-entropy-contraction}.
They distinguish invariant one-time marginals from autonomous Markov histories,
and stationarity from equilibrium. The three-state cycle there shows why
relative-entropy nonincrease need not give convergence. None of these controls
selects Gibbs weighting or supplies a native second law. The physical mixing,
rate and measure requirements remain independent completion tasks.


The paper therefore uses standard equilibrium formulas only after a physical state/measure
gate.  The actual \cmotion{} completion still requires the microscopic closure/open-mode
dynamics, the coarse memory-loss condition if a Markov limit is used, the physical transition
connectivity, and a positive criterion distinguishing thermodynamic equilibrium from other
stationary states.

\section{Radiation as a quantum-statistical control}
\label{sec:radiation}
Radiation tests whether a local-$c$ field theory can recover observed
spectral equilibrium without importing its outcome as a native derivation.
A possible closure-mediated redistribution process is
\begin{equation}
  \boxed{
  \text{radiation mode}
  \;\rightleftarrows\;
  \text{closure internal state}
  \;\rightleftarrows\;
  \text{other radiation modes}.}
  \label{eq:radiation-matter-cycle}
\end{equation}
Such encounters may redistribute frequencies and occupations while preserving
primitive local $c$; they do not select quantum occupation probabilities.

For a standard three-spatial-dimensional free radiation control with linear
dispersion $\omega=c|\mathbf k|$ and $g_{\mathrm{pol}}$ independent polarization
states per wavevector, the density of modes per unit volume and angular-frequency
interval is
\begin{equation}
  g(\omega)\,\dd\omega
  =\frac{g_{\mathrm{pol}}\omega^2}{2\pi^2c^3}\,\dd\omega,
  \qquad g_{\mathrm{pol}}=2\quad\text{(electromagnetic control)}.
  \label{eq:mode-density-control}
\end{equation}
The two electromagnetic polarizations are imported from established physics,
not derived from the primitive \ElEmEr{} propagation premise.  This mode
count is not a count of conserved primitive wave parcels.
\paragraph{The ultraviolet/equipartition gate.}
The same standard mode count exposes why fixed primitive $c$ does not cure
the classical radiation problem. If every independent transverse harmonic
mode instead receives classical equipartition energy $\kb T$, the
electromagnetic ($g_{\mathrm{pol}}=2$) control yields
\begin{equation}
  u_{\omega}^{\mathrm{RJ}}(T)
  =\frac{\kb T\,\omega^2}{\pi^2c^3},\qquad
  \int_0^\infty u_{\omega}^{\mathrm{RJ}}(T)\,\dd\omega=\infty
  \quad(T>0).
  \label{eq:rayleigh-jeans-uv-gate}
\end{equation}
The Rayleigh--Jeans divergence \cite{Rayleigh1900,Jeans1905} comes from unbounded mode count and classical
equipartition, not from variable primitive speed.  Planck's suppression of
high-frequency occupation is an empirical recovery requirement
\cite{Planck1900}.  C-Motion has not derived the requisite occupation,
accessibility or ultraviolet law.  Merely imposing a cutoff or restricting
accessible states would not establish the observed equilibrium distribution.

Suppose additionally that a mode carries the energy increment
\begin{equation}
  \varepsilon_\omega=J_0\omega,
  \label{eq:mode-energy-control}
\end{equation}
and that standard bosonic state counting applies. For a free linear mode at
fixed wavevector, frequency is independent of its occupation/action amplitude.
The assumed increment is therefore compatible with the constant-frequency
exception to the Main Framework's generic nonlinear product-law negative
control, with the vacuum energy subtracted; it is not a product law for a
nonlinear recurrent closure family \cite{CmotionMainFramework2026}.  If the relevant open-mode occupation can
change through a physical reaction network satisfying the full-span hypothesis,
and every exact conserved charge has $q_{\alpha{\rm o}}=0$, the chemical-equilibrium result
\cref{eq:chemical-potential-charge-span} gives zero chemical potential.  The exact phase symmetry stipulated for the Main Framework's recurrent
families is a separate sector assumption: it does not establish an independent
open-mode action charge conserved through closure--radiation conversion.
If a conserved charge of the full coupled system does protect the open-mode
occupation, the zero-chemical-potential control must be reconsidered.
Under the stated additional control assumptions,
\begin{equation}
  \bar n(\omega)
  =\frac{1}{\exp[J_0\omega/(\kb T)]-1},
  \label{eq:planck-occ-control}
\end{equation}
and
\begin{equation}
  u_\omega(T)
  =\frac{J_0\omega^3}{\pi^2c^3}
  \frac{1}{\exp[J_0\omega/(\kb T)]-1}.
  \label{eq:planck-control}
\end{equation}
The displayed spectral density specializes to the electromagnetic control
$g_{\mathrm{pol}}=2$; under otherwise unchanged assumptions a different independent
mode degeneracy would multiply it by $g_{\mathrm{pol}}/2$.
With the empirical Quantum-Foundations identification $J_0=\hbar$, this is the standard
Planck form.  Thermodynamics does not derive that identification, the two-polarization
mode structure, the bosonic occupation law, the required occupation-changing network,
the absence of an independently conserved open-mode action/number, or the physical
mode measure \cite{CmotionQuantumFoundations2026}.

These are recovery controls, not native photon-mode, polarization or thermal
occupation derivations.  No independent carrier frequency, ambient-density
scale or ultraviolet cutoff is calibrated here.

\section{Causal order and the thermodynamic arrow are different questions}
\label{sec:arrows}
The Main Framework uses primitive succession $s$ to express the before--after ordering of
physical change while keeping it distinct from operational clock time
\cite{CmotionMainFramework2026}.  Any entropy derivative $\dd S/\dd s$ below is
schematic along a stipulated ordered physical history; it does not assume that
an observer can read $s$, or that a global equal-$s$ foliation has been derived.
A candidate successor-directed propagation rule says
which later changes can be influenced by a locally generated change.  That is a causal
statement.  It is not yet a statistical statement about entropy.

\begin{proposition}[Successor order does not imply entropy increase]
\label{prop:arrow-separation}
Primitive successor ordering, including a candidate successor-directed propagation
clause for newly generated changes, does not by itself imply
\begin{equation}
  \frac{\dd S}{\dd s}\geq0.
  \label{eq:not-second-law}
\end{equation}
A thermodynamic arrow requires additional reduced-state, statistical, dynamical, and
boundary conditions.
\end{proposition}

Successor order constrains causal influence; a thermodynamic arrow concerns
the behavior of a chosen coarse entropy under physical dynamics and boundary
conditions.  Neither statement is a substitute for the other.

There is nevertheless a useful physical connection that does not collapse the two questions.
An interaction can take compact, easily described incoming organization and distribute its
records into outgoing radiation, closure recoil, internal excitations, and environmental
correlations.  A local subsystem can then look relaxed while the complete extended field
state still carries information about what happened.  Successor-directed propagation helps
move those records away from the subsystem; coarse-graining decides whether they are retained
in the thermodynamic description.  This can contribute to practical macroscopic
irreversibility, but it is not a derivation of entropy increase.

Reversing only the centre-of-momentum motion of a localized system does not
reverse disturbances already propagating outward.  Along the same straight
propagation path, an original outgoing disturbance and a later source-generated
change travelling at the same primitive speed $c$ retain their separation:
the later change cannot overtake the earlier one.  This is a same-path statement,
not a general no-overtaking theorem for different admissible routes.  Reversing
the complete history would require the corresponding distant field structure
and correlations to form a precisely coordinated incoming state.

The Main Framework treats the resulting causal-reachability question as open: a formal
time-reversed field configuration need not be physically reachable from the original
history merely because it satisfies the same mathematical equations.  The required
incoming configuration would itself need a lawful causal ancestry.  A future reachability
constraint could restrict which microscopic histories or transition edges are admissible,
but causal non-reachability by itself does not determine the thermodynamic measure, Gibbs
weighting, the stationary distribution, coarse Markov closure, or an entropy theorem.  It may
therefore constrain the support of thermodynamic dynamics without becoming the second law.

The microscopic reversal/reachability problem is thus separate from selecting
the statistical measure, transition law and conditions for entropy increase.

\section{What remains to complete the thermodynamic bridge}
\label{sec:open}
The present construction already explains why fixed primitive local-$c$ propagation is
compatible with thermodynamics: macroscopic state variables live on the organization and
statistical description of field/closure degrees of freedom rather than on the primitive
speed itself.  A microscopic completion must now supply the parts that the present paper
has deliberately left open.

First, the theory must complete the thermodynamic state-space/measure construction.  The
structural split in \cref{eq:measure-layers} shows that this means more than one task:
the physical support and macro-map, the internal/correlation weighting, the physically
relevant invariant measure, and the equilibrium-selection weighting remain unresolved even
though some closure centre-of-momentum kinematic reference geometry is conditionally constrained.  Second,
the theory must derive the actual closure/open-mode transition map or generator and rates.
It must then establish which kinematically allowed transitions are physically present, whether
unresolved variables lose enough memory for an autonomous coarse closure, and what positive
criterion distinguishes thermodynamic equilibrium from another stationary state.  Third, a
thermal theory of matter requires the closure centre-of-momentum dispersion, the
thermal-activity gates in \cref{sec:thermal-activity-gate}, the kinetic/action
structure and internal excitation spectrum, degeneracies, and occupation-changing
channels.  From the same dynamics
one would then seek fluctuation--dissipation relations and transport coefficients such as
diffusion, viscosity, conductivity, and heat transport.

Quantum radiation adds a further layer: the microscopic occupation law and the observed
Planck spectrum must be recovered without treating the downstream control formulas as a
microscopic derivation.  Finally, a complete thermodynamic arrow requires the statistical,
dynamical, and boundary antecedents of a second law, while third-law and cosmological
thermal questions remain later extensions.

No independent thermodynamic \cmotion{} scale is introduced to bypass these problems.
Any framework-scale input must obey the common one-scale policy of the Main Framework;
this paper does not calibrate a separate carrier frequency, wavelength, density scale,
cutoff, or closure thermal scale.

\section{Conclusion}
\label{sec:conclusion}
The primitive local-$c$ constraint is compatible with thermodynamics because
thermal variation concerns collective closure motion, internal excitations,
open-mode spectra, correlations and the statistical weights of admissible
physical configurations, not a distribution of primitive speeds.

The field-first ontology rules out treating arbitrary ray, crest or preparation
labels as independent entropy constituents.  The M1--M4 decomposition isolates
the physical state-space, kinematic reference, invariant-measure and equilibrium
weighting questions.  Conditional effective Lorentz symmetry constrains a
closure centre-of-momentum reference measure, but neither this measure nor
invariance itself selects thermodynamic weighting.  Recurrent closure geometry
also does not determine its entropy without a physical sector measure.

The T0--T5 hierarchy distinguishes dephasing, scattering, energy exchange,
occupation change and sustained mixing.  Closure-mediated interactions are a
candidate way of realizing these processes; their native generator, rates,
connectivity, coarse memory loss and equilibrium selector remain unresolved.
A deterministic complete-state evolution can yield stochastic coarse transitions,
but a fixed one-step kernel does not establish Markovian multi-step histories.
Stationarity, even with preparation loss, does not by itself identify thermal
equilibrium.

Energy, action and entropy remain distinct accounts.  Conventional heat/work,
kinetic-gas and radiation formulas apply only after their conservation,
measure, dispersion, mode-counting and occupation assumptions are supplied.
In particular, fixed primitive $c$ does not remove the classical ultraviolet
radiation problem, and the Planck distribution is a downstream control rather
than a native derivation.  Causal succession and the thermodynamic arrow are
likewise different: outgoing correlations may disperse records without proving
a second-law inequality.

The outstanding task is to derive from the microscopic \ElEmEr{} dynamics the
physical reduction and weighting, active closure spectrum, allowed transitions
and rates, justified coarse evolution and equilibrium selection, quantum
occupation law, transport and robust macroscopic irreversibility.  These remain
Open/Conditional/Deferred at the scopes recorded in the formal register.  The
paper's result is a constrained thermodynamic research framework, not a
completed microscopic derivation of equilibrium or entropy production.

\newpage

\appendix

\section{Formal nativity and claim register}
\label{sec:register}
This register records scientific status and nativity independently.  Standard
controls, model assumptions and inherited results are identified separately;
no native Open or Conditional result is promoted.

\paragraph{How to read the register.}
Scientific status follows the Derivation Ledger's vocabulary:
\emph{Structural} denotes exact or banked organizational/accounting statements,
not a microscopic law; \emph{Conditional} denotes conclusions from explicit
premises or model classes not themselves derived; \emph{Negative} excludes an
identification or mechanism only within its stated scope. \emph{Conjectural}
denotes a motivated but underived interpretation; \emph{Open} is unresolved,
and \emph{Open (sharp)} isolates a missing bridge rather than greater confidence.
\emph{Deferred} places work outside the present publication critical path
without resolving its scientific status.

The independent provenance shorthand is read locally as follows:
\emph{N1} identifies direct use of the stated C-Motion physical/organizational
premises, not a derivation of those premises; \emph{N2} identifies calculations
or applications in an explicitly assumed C-Motion-motivated model class;
\emph{C} denotes a downstream standard or externally supplied mathematical or
physical control; \emph{O} marks an open native input or completion.
The notation \emph{N2 $\to$ C} means that a conditional C-Motion-motivated
setup is followed by standard mathematics, not a scientific-status promotion.
Provenance codes do not establish native microscopic dynamics.

\begin{center}
\footnotesize
\renewcommand{\arraystretch}{1.20}
\begin{longtable}{@{}P{3.5cm}P{1.7cm}P{2.5cm}P{5.6cm}@{}}
\toprule
Result & Nativity & Scientific status & Scope / limitations \\
\midrule
\endfirsthead
\toprule
Result & Nativity & Scientific status & Scope / limitations \\
\midrule
\endhead
Field-first fine physical state
& \natNOne{} & \statusStructural{}
& The physical microstate is a complete field/closure configuration, not an immutable parcel list. \\
No automatic primitive crest/ray/packet-count entropy
& \natNOne{} & \statusNegative{}
& Primitive descriptive labels do not supply a physical entropy count without a derived reduced state space and measure. \\
Open/closed/extended organizational distinction
& \natNOne{} & \statusStructural{}
& Thermodynamic reduction may depend on organization; the full reduced state space is not yet derived. \\
Closure thermal-activity diagnostic
& \natNTwo{} / \natControl{} & \statusStructural{} / \statusOpen{}
& A physical coordinate counts as thermally active only with redundancy and sector guards, finite excitation, dynamics and interaction connectivity; the native closure spectrum remains Open. \\
Quadratic closure-spectrum test
& \natControl{} & \statusConditional{}
& In a supplied second-order quadratic model $\det(H-\Omega^2G)=0$; the Hessian alone does not fix frequencies, and no native kinetic/action metric is derived. \\
Recurrent organization directly fixes closure entropy
& \natControl{} & \statusNegative{}
& Codimension or zero ambient measure is not an entropy calculation; a justified induced/conditional sector measure and fixed reference normalization are required. \\
Physical reduced-state / macro-map / measure gate
& \natNTwo{} & \statusConditional{} / \statusOpen{}
& Generic statistical requirements are applied only after \cmotion{} ontology constrains the admissible state space. \\
Measure problem decomposition (M1--M4)
& \natNTwo{} / \natControl{} & \statusStructural{} / \statusOpen{}
& Separates physical support, kinematic reference measure, dynamically invariant measure, and thermodynamic equilibrium weighting.  Some COM kinematic geometry is conditionally constrained; the full physical/equilibrium measure is not. \\
Invariant measure uniquely selects thermodynamic weighting
& \natControl{} & \statusNegative{}
& Examples with fixed points, periodic orbits or mixtures show generic nonselection; uniqueness of an invariant probability in some systems does not itself give a physical equilibrium interpretation. \\
Conditional closure-COM kinematic reference measure
& \natNTwo{} $\to$ \natControl{} & \statusConditional{}
& In the already-Conditional effective massive/Lorentz sector, the positive mass shell carries $\dd^3P/p^0$ up to normalization.  Densities relative to this measure and to $\dd^3P$ differ by the $p^0$ factor; neither convention alone selects equilibrium. \\
Chemical-potential structure from conserved charges
& \natControl{}; application \natNTwo{} & \statusConditional{}
& Allowed reaction directions imply $\sum_i\nu_i^{(r)}\mu_i=0$; under $\operatorname{span}\{\nu^{(r)}\}=\ker q$, $\mu\in\operatorname{im}q^{\mathsf T}$. If every $q_{\alpha{\rm o}}=0$, then $\mu_{\rm o}=0$; occupation change alone is insufficient. \\
Deterministic-fine to coarse-stochastic bridge
& \natNTwo{} $\to$ \natControl{} & \statusConditional{}
& In a supplied compatible complete-state/evolution representation, a normalized fine-state probability measure and measurable positive-weight cells define a row-stochastic one-step kernel; neither global succession alignment nor autonomous multi-step coarse dynamics follows. \\
Coarse Markovianity is automatic after projection
& \natControl{} & \statusNegative{}
& A fixed one-step matrix does not imply autonomous multi-step Markov histories; exact lumpability or justified approximate memory loss must be established separately. \\
Stationarity/preparation independence identifies equilibrium
& \natControl{} & \statusNegative{}
& $\pi W=\pi$ and even $P_n\to\pi$ do not by themselves identify thermodynamic equilibrium; stationary circulating currents are possible. \\
Connectivity/balance does not select Gibbs
& \natControl{} & \statusNegative{}
& Abstract connectivity or balance structure can support different stationary weights; the physical measure/rate structure must supply the equilibrium energy dependence. \\

Two-bath equal temperature and stability
& \natControl{} & \statusConditional{}
& Supplied weak-coupling multiplicity product and differentiable entropies give equal temperature at an interior stationary partition; positive heat capacities make it a local maximum, not a global uniqueness theorem. \\
Data processing and contraction $\neq$ second law
& \natControl{} & \statusConditional{} / \statusNegative{}
& For an autonomous Markov kernel with stationary $\pi$ and defined finite relative entropy, $D(PW\Vert\pi)\leq D(P\Vert\pi)$; neither strict contraction, convergence, Gibbs selection nor a native second law follows. \\
Kinetic arrest $\neq$ third law
& \natControl{} & \statusNegative{}
& Slow or absent transitions do not establish equilibrium ground-state entropy or a native third law; disconnected sectors need a justified physical counting rule. \\
Coarse-graining correlation redistribution
& \natControl{}; application \natNTwo{} & \statusConditional{}
& \Cref{eq:correlation-redistribution} interprets information discarded by a supplied macro-description; it is not a native monotonic-entropy theorem. \\

Boltzmann/Gibbs entropy forms
& \natControl{} & \statusConditional{}
& Standard formulas given a reduced state space, fixed dimensionless reference-measure convention and normalized probability density; native physical weighting and state-cell normalization remain open. \\
Primitive local $c$ does not determine temperature
& \natNOne{} & \statusNegative{}
& Thermodynamic variation resides in energy distribution, organization, correlations, composite motion, and occupation, not primitive speed. \\
$1/T=\left(\partial S/\partial U\right)_{\{\mathcal X\}}$
& \natControl{} & \statusStructural{}
& Introduced after the native fixed-$c$ restriction; not derived from $c$. \\
Operational thermometer agreement
& \natControl{} & \statusOpen{}
& Independent weakly coupled probes must agree in a specified contact/equilibrium regime; no native thermalizer or thermometer response is supplied. \\
T0--T5 generic hierarchy
& \natControl{} & \statusStructural{}
& Diagnostic taxonomy; not by itself native evidence. \\
Assignment of open/closure processes to T0--T5
& \natNTwo{} & \statusConditional{} / \statusOpen{}
& Native process classes constrain which dynamical jobs can be performed. \\
Free propagation/dephasing gives full thermalization
& \natControl{} & \statusNegative{}
& T0/T1 can relax selected coarse observables but do not establish generic T2--T5; T5 additionally does not collapse connectivity, coarse closure, stationarity, and equilibrium into one condition. \\
Closure/matter-mediated equilibration route
& \natNTwo{} & \statusConditional{}
& Process-sufficiency class only; transition law/rates, physical connectivity, coarse memory-loss/closure, and equilibrium selection remain Open. \\
Free-propagation class / interaction-opportunity thermalizer architecture
& \natNTwo{} candidate & \statusConjectural{} / \statusOpen{}
& Possible microscopic realization; not shown to supply T1--T5, detailed balance, or binding. \\
One physical energy ledger; energy/action/entropy distinct
& \natNTwo{} & \statusStructural{} / \statusConditional{}
& A common conserved energy functional and complete-system exchange balance are conditional inputs, not a derived native field law; $J_0$ is not entropy or temperature. \\
First-law algebra $\dd U=\delta Q+\delta W$
& \natControl{} & \statusStructural{}
& Exact two-term accounting identity within the conditionally supplied conserved energy ledger, specified boundary and complete allocation to heat/work (or additional explicit terms); Structural classifies the algebra, not native energy conservation. \\
Radiated energy is automatically heat
& \natControl{} & \statusNegative{}
& Radiation can be coherent/work-capable or correlated; ``heat'' depends on subsystem and macro-description. \\
Canonical weighting
& \natControl{} & \statusConditional{}
& Requires specified $H$, measure and equilibrium antecedents. \\
Standard gas distribution controls
& \natControl{} & \statusConditional{}
& Downstream of assumed timelike closure dispersion and a supplied rest-frame equilibrium prescription; the Maxwell--Boltzmann limit additionally requires a nondegenerate classical and $\kb T\ll Mc^2$ thermal regime. \\
Detailed/global balance controls
& \natControl{} & \statusStructural{} / \statusOpen{}
& A normalized invariant fine measure yields the one-step identity $\pi W=\pi$ even without autonomous coarse Markov closure; physical multi-step propagation requires additional closure assumptions.  Pairwise detailed balance and thermodynamic equilibrium remain separate questions. \\
Classical equipartition cures ultraviolet radiation
& \natControl{} & \statusNegative{}
& For imported two-polarization, unbounded linear-mode counting, $u_\omega^{\rm RJ}=\kb T\omega^2/(\pi^2c^3)$ diverges on integration; a native occupation/UV/accessibility mechanism remains Open. \\
Planck/radiation formulas
& \natControl{} & \statusConditional{}
& Requires imported two-transverse-polarization electromagnetic mode count $g_{\mathrm{pol}}=2$, $J_0\omega$, bosonic occupation statistics and the normalization status in Quantum Foundations; not a native ElEmEr photon-mode derivation. \\
Causal/successor orientation is not the thermodynamic arrow
& \natNOne{} & \statusStructural{}
& Primitive succession does not imply $\dd S/\dd s\ge0$; second-law antecedents remain Open. \\
Native reduced state space and physical thermodynamic measure
& \natOpen{} & \statusSharpOpen{}
& Physical support, internal/correlation weighting, selected invariant measure, and equilibrium weighting remain unresolved; symmetry constrains only some conditional kinematic reference measures. \\
Native transition dynamics, coarse closure, connectivity, rates, and equilibrium selection
& \natOpen{} & \statusOpen{}
& Distinct unresolved requirements: actual physical edges/rates, loss of dynamically relevant hidden memory, connected exploration, and the positive criterion selecting thermodynamic equilibrium rather than another stationary state. \\
Closure internal spectrum and transport coefficients
& \natOpen{} & \statusOpen{} / \statusDeferred{}
& Genuine completion work; may remain unresolved in this explicitly work-in-progress manuscript. \\
Native quantum occupation law / Planck recovery
& \natOpen{} & \statusOpen{} / \statusDeferred{}
& Must not be inferred from primitive $c$, recurrence, or downstream controls. \\
Arrow, third-law and cosmological completion
& \natOpen{} & \statusOpen{} / \statusDeferred{}
& Explicitly outside the scope required for the present work-in-progress manuscript. \\
\bottomrule
\end{longtable}
\normalsize
\end{center}

\subsection{Status note}
The register separates scoped standard controls, conditional applications and negative results;
no pre-existing Open or Conditional native thermodynamic result is promoted.  The complete
physical state space/support, internal/correlation measure, physical invariant/equilibrium
weighting, actual transition map and rates, coarse memory-loss/closure mechanism, physical
connectivity, positive equilibrium-selection criterion, closure spectrum, transport
coefficients, quantum occupation law, Planck recovery, rigorous second-law theorem,
microscopic irreversibility account, low-entropy boundary, third-law completion, and
cosmological thermal history remain unresolved.  The closure/matter-mediated equilibration
picture remains Conditional, and the free-propagation/interaction-opportunity realization remains
Conjectural/Open.  The status of $J_0=\hbar$ remains inherited from Quantum Foundations
rather than rederived here.

\newpage

\printbibliography[title={References}]

\end{document}
